%=========================================================
% C++ Chapter 12 Lab Handout (Verbatim Korean-safe)
% Topic: File I/O (Text/Binary, Modes, State, Random Access)
% Compile with XeLaTeX
%=========================================================
\documentclass[11pt,a4paper]{article}

\usepackage[margin=1in]{geometry}
\usepackage{fontspec}
\usepackage{kotex}
\usepackage{hyperref}
\usepackage{enumitem}
\usepackage{fancyhdr}
\usepackage{titlesec}
\usepackage{xcolor}
\usepackage{tcolorbox}
\usepackage{fvextra}

% ---------- Fonts ----------
% If D2Coding isn't available, try: NanumGothicCoding, Noto Sans Mono CJK KR
\setmonofont{D2Coding}[Scale=MatchLowercase]

% ---------- Colors ----------
\definecolor{CodeBg}{RGB}{248,248,248}
\definecolor{CodeFrame}{RGB}{220,220,220}
\definecolor{Accent}{RGB}{40,80,160}

% ---------- Header ----------
\pagestyle{fancy}
\fancyhf{}
\lhead{\textbf{C++ 12장 실습}}
\rhead{\textbf{C++ 파일 입출력}}
\cfoot{\thepage}

\titleformat{\section}{\Large\bfseries\color{Accent}}{\thesection}{0.6em}{}
\titleformat{\subsection}{\large\bfseries}{\thesubsection}{0.6em}{}

% ---------- Boxes ----------
\newtcolorbox{GoalBox}{
  colback=CodeBg,
  colframe=CodeFrame,
  arc=3pt,
  boxrule=0.8pt,
  left=10pt,right=10pt,top=8pt,bottom=8pt
}
\newcommand{\filetag}[1]{\texttt{\color{Accent}#1}}

% ---------- Code block environment (Unicode-safe) ----------
% IMPORTANT: do NOT set commandchars, so C++ braces {} print correctly.
\DefineVerbatimEnvironment{cppcode}{Verbatim}{
  fontsize=\small,
  frame=single,
  framerule=0.6pt,
  rulecolor=\color{CodeFrame},
  numbers=left,
  numbersep=8pt,
  baselinestretch=1,
  breaklines=true,
  breaksymbolleft=,
  breaksymbolright=,
  xleftmargin=1.5em
}

\begin{document}

\begin{center}
  {\LARGE \textbf{C++ 12장 실습: 파일 입출력}}\\
  \vspace{6pt}
  {\large (text vs binary / fstream / file modes / read-write / stream state / random access)}
\end{center}

\vspace{8pt}

\section*{학습 목표(슬라이드 기반)}
\begin{GoalBox}
\begin{enumerate}[leftmargin=18pt]
  \item 텍스트 파일과 바이너리 파일의 차이를 안다.
  \item C++ 표준 파일 입출력 라이브러리(\texttt{ifstream/ofstream/fstream})를 안다.
  \item \texttt{<<, >>}로 간단히 텍스트 파일을 읽고 쓰는 방법을 안다.
  \item 파일 모드(file mode)를 이해한다.
  \item 텍스트 I/O 모드로 파일을 읽고 쓰는 방법을 안다.
  \item 바이너리 I/O 모드로 파일을 읽고 쓰는 방법을 안다.
  \item 텍스트 I/O와 바이너리 I/O 차이를 이해한다.
  \item 파일 입출력 스트림의 상태를 검사하는 방법을 안다.
  \item 임의 접근(random access)으로 파일을 읽고 쓰는 방법을 안다.
\end{enumerate}
\end{GoalBox}

\section*{실습 운영 규칙}
\begin{itemize}[leftmargin=18pt]
  \item 모든 예제는 \texttt{main()} 포함, \textbf{독립 실행}입니다.
  \item 윈도우 경로(\texttt{c:\textbackslash temp\textbackslash ...}) 대신, 수업에서는 \textbf{상대 경로}(\texttt{./data/...})를 사용합니다.
  \item 바이너리 예제는 이미지 파일이 없어도 실습할 수 있도록 \textbf{일반 파일 복사} 버전도 제공합니다.
  \item 텍스트 I/O에서 \texttt{\textbackslash n} 처리(윈도우의 \texttt{\textbackslash r\textbackslash n}) 차이를 꼭 관찰하세요.
\end{itemize}

\tableofcontents
\newpage

%=========================================================
\section{텍스트 vs 바이너리: 핵심 차이(Enter, \textbackslash r\textbackslash n)}
\begin{GoalBox}
\textbf{핵심}
\begin{itemize}[leftmargin=18pt]
  \item 텍스트 파일: ``문자''로 해석 가능한 데이터 중심(ASCII/Unicode). % p.4
  \item 바이너리 파일: 문자로 해석되지 않는 바이트가 포함(이미지/오디오/문서/실행파일). % p.8~10
  \item 윈도우 텍스트 파일은 줄바꿈을 \texttt{\textbackslash r\textbackslash n} 두 바이트로 저장하지만, 텍스트 I/O는 보통 \texttt{\textbackslash n}로 취급. % p.7, p.21
  \item 바이너리 I/O는 바이트를 있는 그대로 읽고 씀(\texttt{ios::binary}). % p.28, p.35
\end{itemize}
\end{GoalBox}

%=========================================================
\section{파일 스트림 클래스: ifstream/ofstream/fstream}
\begin{GoalBox}
\textbf{핵심}
\begin{itemize}[leftmargin=18pt]
  \item \texttt{<fstream>} 헤더 + \texttt{std} 네임스페이스. % p.14
  \item \texttt{ifstream}: 읽기, \texttt{ofstream}: 쓰기, \texttt{fstream}: 읽기+쓰기. % p.11
  \item 파일과 프로그램을 ``스트림''으로 연결하여 \texttt{>>}, \texttt{get/read}, \texttt{<<}, \texttt{put/write} 사용. % p.13
\end{itemize}
\end{GoalBox}

%=========================================================
\section{예제 12-1: 키보드 입력을 텍스트 파일로 저장(<<)}
\begin{GoalBox}
\textbf{핵심}
\begin{itemize}[leftmargin=18pt]
  \item \texttt{ofstream fout("student.txt");} 처럼 생성과 동시에 open 가능. % p.16~17
  \item \texttt{if(!fout)} 또는 \texttt{if(!fout.is\_open())}로 열기 실패 검사. % p.16~18
  \item 정수도 텍스트로 저장되며(숫자 문자로 변환), 다음에 \texttt{>>}로 다시 읽을 수 있음. % p.17
\end{itemize}
\end{GoalBox}

\noindent\filetag{L1\_write\_student\_text.cpp}
\begin{cppcode}
#include <iostream>
#include <fstream>
#include <string>
using namespace std;

/*
[포인트]
- 교재는 c:\temp\student.txt 같은 절대경로를 사용하지만,
  수업에서는 프로젝트 폴더 기준 상대경로(./data/student.txt)를 권장.
- 폴더가 없으면 파일 생성이 실패할 수 있으니, ./data 폴더를 미리 만들어 두기.
*/

int main() {
    string name, dept;
    int sid;

    cout << "이름>> ";
    cin >> name;
    cout << "학번>> ";
    cin >> sid;
    cout << "학과>> ";
    cin >> dept;

    ofstream fout("./data/student.txt");
    if(!fout) {
        cout << "파일 열기 실패: ./data/student.txt\n";
        return 0;
    }

    fout << name << '\n';
    fout << sid  << '\n';
    fout << dept << '\n';
    fout.close();

    cout << "저장 완료\n";
}
\end{cppcode}

%=========================================================
\section{예제 12-2: 텍스트 파일 읽기(>>)}
\begin{GoalBox}
\textbf{핵심}
\begin{itemize}[leftmargin=18pt]
  \item \texttt{>>}는 공백 기준으로 토큰을 읽는다(공백 포함 라인은 getline 사용). % 11장과 연결
  \item 파일 열기 실패는 ``경로 오류/파일 없음''이 흔한 원인. % p.18
\end{itemize}
\end{GoalBox}

\noindent\filetag{L2\_read\_student\_text.cpp}
\begin{cppcode}
#include <iostream>
#include <fstream>
#include <string>
using namespace std;

int main() {
    ifstream fin("./data/student.txt");
    if(!fin) {
        cout << "파일 열기 실패: ./data/student.txt\n";
        return 0;
    }

    string name, dept;
    int sid;

    fin >> name;
    fin >> sid;
    fin >> dept;

    cout << name << "\n" << sid << "\n" << dept << "\n";
    fin.close();
}
\end{cppcode}

%=========================================================
\section{파일 모드(file mode)와 open()}
\begin{GoalBox}
\textbf{핵심 모드}
\begin{itemize}[leftmargin=18pt]
  \item \texttt{ios::in}: 읽기, \texttt{ios::out}: 쓰기
  \item \texttt{ios::app}: 파일 끝에 덧붙이기(append) % p.20
  \item \texttt{ios::trunc}: 기존 내용을 지우고 새로 쓰기(보통 ofstream 기본 동작) % p.19
  \item \texttt{ios::binary}: 바이너리 I/O(텍스트 변환 금지) % p.20, p.28
\end{itemize}
\end{GoalBox}

\noindent\filetag{L3\_append\_mode.cpp}
\begin{cppcode}
#include <iostream>
#include <fstream>
using namespace std;

/*
[포인트]
- ios::app : 항상 파일 끝에 쓴다.
- ios::out | ios::app 로 열고 한 줄을 추가해 보자.
*/

int main() {
    ofstream fout("./data/student.txt", ios::out | ios::app);
    if(!fout) {
        cout << "파일 열기 실패\n";
        return 0;
    }
    fout << "tel:010-0000-0000\n";
    fout.close();
    cout << "append 완료\n";
}
\end{cppcode}

%=========================================================
\section{예제 12-3: get()으로 텍스트 파일을 바이트 단위 읽기(텍스트 I/O)}
\begin{GoalBox}
\textbf{핵심}
\begin{itemize}[leftmargin=18pt]
  \item \texttt{fin.get()}은 EOF(-1)까지 1바이트씩 읽는다. % p.21~23
  \item \textbf{텍스트 I/O 모드}에서는 윈도우의 \texttt{\textbackslash r\textbackslash n}이 읽힐 때 \texttt{\textbackslash n}만 리턴되는 효과가 나타날 수 있음(교재 예: 바이트 수가 줄어듦). % p.21
\end{itemize}
\end{GoalBox}

\noindent\filetag{L4\_get\_text\_count.cpp}
\begin{cppcode}
#include <iostream>
#include <fstream>
using namespace std;

/*
[실습]
- ./data/sample.txt 를 미리 만들어 두고 실행하세요.
- (윈도우) 텍스트 모드에서 \r\n이 \n처럼 처리되면 카운트가 달라질 수 있음.
*/

int main() {
    const char* file = "./data/sample.txt";
    ifstream fin(file); // 디폴트: ios::in (텍스트)
    if(!fin) {
        cout << "열기 오류: " << file << "\n";
        return 0;
    }

    int count = 0;
    int c;
    while((c = fin.get()) != EOF) {
        cout.put((char)c);
        count++;
    }
    cout << "\n읽은 바이트 수(텍스트 모드 count) = " << count << "\n";
    fin.close();
}
\end{cppcode}

%=========================================================
\section{예제 12-4: 텍스트 파일 연결(한 파일 내용을 다른 파일 끝에 붙이기)}
\begin{GoalBox}
\textbf{핵심}
\begin{itemize}[leftmargin=18pt]
  \item \texttt{fstream}은 읽기/쓰기를 모두 지원. % p.24
  \item 입력 파일을 \texttt{get()}으로 읽고 출력 파일에 \texttt{put()}으로 기록.
  \item 출력 파일은 \texttt{ios::app}로 열어 ``끝에 덧붙이기''.
\end{itemize}
\end{GoalBox}

\noindent\filetag{L5\_concat\_text\_files.cpp}
\begin{cppcode}
#include <iostream>
#include <fstream>
using namespace std;

int main() {
    const char* firstFile  = "./data/student.txt";
    const char* secondFile = "./data/sample.txt";

    fstream fout(firstFile, ios::out | ios::app);
    if(!fout) { cout << "열기 오류: " << firstFile << "\n"; return 0; }

    fstream fin(secondFile, ios::in);
    if(!fin) { cout << "열기 오류: " << secondFile << "\n"; return 0; }

    int c;
    while((c = fin.get()) != EOF) {
        fout.put((char)c);
    }

    fin.close();
    fout.close();
    cout << "연결(append) 완료\n";
}
\end{cppcode}

%=========================================================
\section{라인 단위 읽기: getline 2가지 방식(예제 12-5, 12-6)}
\begin{GoalBox}
\textbf{핵심}
\begin{itemize}[leftmargin=18pt]
  \item \texttt{fin.getline(char*, n)}: C 스타일 버퍼로 라인 읽기. % p.25~26
  \item \texttt{getline(fin, string)}: 전역 함수로 string에 라인 읽기(더 안전). % p.25~27
\end{itemize}
\end{GoalBox}

\noindent\filetag{L6\_getline\_charbuf.cpp}
\begin{cppcode}
#include <iostream>
#include <fstream>
using namespace std;

int main() {
    ifstream fin("./data/sample.txt");
    if(!fin) { cout << "열기 실패\n"; return 0; }

    char buf[81]; // 한 줄 최대 80자 가정
    while(fin.getline(buf, 81)) {
        cout << buf << "\n";
    }
    fin.close();
}
\end{cppcode}

\noindent\filetag{L6b\_getline\_string.cpp}
\begin{cppcode}
#include <iostream>
#include <fstream>
#include <string>
#include <vector>
using namespace std;

static void fileRead(vector<string>& v, ifstream& fin) {
    string line;
    while(getline(fin, line)) {
        v.push_back(line);
    }
}

static void search(const vector<string>& v, const string& word) {
    for (int i = 0; i < (int)v.size(); i++) {
        if (v[i].find(word) != string::npos)
            cout << v[i] << "\n";
    }
}

int main() {
    ifstream fin("./data/words.txt");
    if(!fin) { cout << "words.txt 열기 실패(./data/words.txt)\n"; return 0; }

    vector<string> wordVector;
    fileRead(wordVector, fin);
    fin.close();

    cout << "words.txt 파일을 읽었습니다.\n";
    while(true) {
        cout << "검색할 단어를 입력하세요 >> ";
        string word;
        getline(cin, word);
        if(word == "exit") break;
        search(wordVector, word);
    }
    cout << "프로그램을 종료합니다.\n";
}
\end{cppcode}

%=========================================================
\section{바이너리 I/O: get/put로 파일 복사(예제 12-7)}
\begin{GoalBox}
\textbf{핵심}
\begin{itemize}[leftmargin=18pt]
  \item \texttt{ios::binary}를 빼면(특히 윈도우) 줄바꿈 변환 등으로 \textbf{원본과 동일한 바이트 복사}가 깨질 수 있음. % p.29, p.35
  \item 바이너리 파일뿐 아니라 텍스트 파일도 바이너리 모드로 복사 가능.
\end{itemize}
\end{GoalBox}

\noindent\filetag{L7\_copy\_binary\_getput.cpp}
\begin{cppcode}
#include <iostream>
#include <fstream>
using namespace std;

int main() {
    const char* src  = "./data/src.bin";   // 아무 파일이나 OK (텍스트/바이너리)
    const char* dest = "./data/dest.bin";

    ifstream fsrc(src, ios::in | ios::binary);
    if(!fsrc) { cout << "열기 오류: " << src << "\n"; return 0; }

    ofstream fdest(dest, ios::out | ios::binary);
    if(!fdest) { cout << "열기 오류: " << dest << "\n"; return 0; }

    int c;
    while((c = fsrc.get()) != EOF) {
        fdest.put((char)c);
    }

    fsrc.close();
    fdest.close();
    cout << "복사 완료: " << src << " -> " << dest << "\n";
}
\end{cppcode}

%=========================================================
\section{read/write: 블록 단위 복사(예제 12-9)}
\begin{GoalBox}
\textbf{핵심}
\begin{itemize}[leftmargin=18pt]
  \item \texttt{read(buf, N)}는 최대 N바이트를 읽고, 실제 읽은 바이트는 \texttt{gcount()}로 확인. % p.30~32
  \item \texttt{eof()}는 ``읽기 시도 이후''에 true가 될 수 있어, 루프 설계를 주의.
\end{itemize}
\end{GoalBox}

\noindent\filetag{L8\_copy\_binary\_readwrite.cpp}
\begin{cppcode}
#include <iostream>
#include <fstream>
using namespace std;

int main() {
    const char* src  = "./data/src.bin";
    const char* dest = "./data/dest_readwrite.bin";

    ifstream fsrc(src, ios::in | ios::binary);
    if(!fsrc) { cout << "열기 오류: " << src << "\n"; return 0; }

    ofstream fdest(dest, ios::out | ios::binary);
    if(!fdest) { cout << "열기 오류: " << dest << "\n"; return 0; }

    char buf[1024];
    while(true) {
        fsrc.read(buf, sizeof(buf));
        streamsize n = fsrc.gcount(); // 실제 읽은 바이트
        if(n <= 0) break;
        fdest.write(buf, n);
    }

    fsrc.close();
    fdest.close();
    cout << "블록 복사 완료\n";
}
\end{cppcode}

%=========================================================
\section{바이너리 데이터 저장/복원: int 배열 + double (예제 12-10)}
\begin{GoalBox}
\textbf{핵심}
\begin{itemize}[leftmargin=18pt]
  \item 숫자를 텍스트로 저장하면 ``문자''로 변환되어 저장(사람이 읽기 쉬움).
  \item 바이너리로 저장하면 메모리의 바이트가 그대로 기록(빠르고 크기 작음) % p.33
  \item 단, 바이너리는 엔디안/타입 크기/구조체 패딩에 영향을 받을 수 있어 \textbf{호환성 이슈}가 존재.
\end{itemize}
\end{GoalBox}

\noindent\filetag{L9\_save\_load\_binary\_numbers.cpp}
\begin{cppcode}
#include <iostream>
#include <fstream>
using namespace std;

int main() {
    const char* file = "./data/data.dat";

    int n[10] = {0,1,2,3,4,5,6,7,8,9};
    double d = 3.15;

    // 저장
    ofstream fout(file, ios::out | ios::binary);
    if(!fout) { cout << "열기 오류\n"; return 0; }
    fout.write((char*)n, sizeof(n));
    fout.write((char*)(&d), sizeof(d));
    fout.close();

    // 값 변경(복원 확인용)
    for(int i=0;i<10;i++) n[i] = 99;
    d = 8.15;

    // 복원
    ifstream fin(file, ios::in | ios::binary);
    if(!fin) { cout << "열기 오류\n"; return 0; }
    fin.read((char*)n, sizeof(n));
    fin.read((char*)(&d), sizeof(d));
    fin.close();

    for(int i=0;i<10;i++) cout << n[i] << ' ';
    cout << "\n" << d << "\n";
}
\end{cppcode}

%=========================================================
\section{스트림 상태 검사: good/fail/bad/eof}
\begin{GoalBox}
\textbf{핵심}
\begin{itemize}[leftmargin=18pt]
  \item \texttt{if(!fin)}는 \texttt{fail()}을 포함한 상태 이상을 의미하는 빠른 체크.
  \item \texttt{eof()}는 ``EOF를 이미 만났는가''이며, 보통 \textbf{읽기 시도 이후}에 true가 된다. % p.22, p.31
  \item 안전한 패턴: \texttt{while(getline(fin,line))} / \texttt{while(fin.read(...), fin.gcount()>0)}.
\end{itemize}
\end{GoalBox}

\noindent\filetag{L10\_stream\_state.cpp}
\begin{cppcode}
#include <iostream>
#include <fstream>
#include <string>
using namespace std;

int main() {
    ifstream fin("./data/sample.txt");
    if(!fin) { cout << "open fail\n"; return 0; }

    string line;
    while(true) {
        if(!getline(fin, line)) {
            if(fin.eof()) cout << "EOF 도달\n";
            else cout << "읽기 실패(fail/bad)\n";
            break;
        }
        cout << line << "\n";
    }
}
\end{cppcode}

%=========================================================
\section{임의 접근(Random Access): seekg/tellg, seekp/tellp}
\begin{GoalBox}
\textbf{핵심}
\begin{itemize}[leftmargin=18pt]
  \item \texttt{tellg()} / \texttt{tellp()} : 현재 파일 포인터 위치(바이트 오프셋) % 11장의 개념 확장
  \item \texttt{seekg(pos)} / \texttt{seekp(pos)} : 파일 포인터 이동
  \item 텍스트 파일은 OS의 줄바꿈 변환 때문에 정확한 바이트 위치가 직관적이지 않을 수 있어,
        \textbf{임의 접근은 바이너리 모드에서 더 안전}.
\end{itemize}
\end{GoalBox}

\noindent\filetag{L11\_random\_access\_binary.cpp}
\begin{cppcode}
#include <iostream>
#include <fstream>
using namespace std;

/*
[실습]
- data.dat(L9)이 생성된 뒤에 실행하면 좋습니다.
- 첫 int 10개(40바이트) + double(8바이트) 구조.
- 특정 위치로 이동하여 일부만 읽어봅니다.
*/

int main() {
    const char* file = "./data/data.dat";
    ifstream fin(file, ios::in | ios::binary);
    if(!fin) { cout << "open fail\n"; return 0; }

    // 5번째 int(0-based index 4) 읽기: offset = 4 * sizeof(int)
    fin.seekg(4 * (int)sizeof(int), ios::beg);
    int x;
    fin.read((char*)&x, sizeof(x));
    cout << "5번째 int = " << x << "\n";

    // double 읽기: offset = sizeof(int)*10
    fin.seekg(10 * (int)sizeof(int), ios::beg);
    double d;
    fin.read((char*)&d, sizeof(d));
    cout << "double = " << d << "\n";

    fin.close();
}
\end{cppcode}

\section*{마무리 체크리스트}
\begin{GoalBox}
\begin{itemize}[leftmargin=18pt]
  \item 텍스트/바이너리 파일의 차이와 줄바꿈(\texttt{\textbackslash r\textbackslash n}) 이슈를 설명할 수 있는가?
  \item \texttt{ifstream/ofstream/fstream}의 역할 차이를 말할 수 있는가?
  \item 파일 모드(\texttt{ios::in/out/app/trunc/binary})를 상황에 맞게 선택할 수 있는가?
  \item get/put(1바이트)과 read/write(블록)의 차이를 설명할 수 있는가?
  \item \texttt{eof()}를 루프 조건으로 쓰는 것이 위험한 이유를 말할 수 있는가?
  \item 바이너리 저장의 장점/단점(호환성)을 말할 수 있는가?
  \item \texttt{tellg/seekg}로 임의 접근을 구현할 수 있는가?
\end{itemize}
\end{GoalBox}

\end{document}
